%%
%% paper78_whitehead_aspharizitaet_de.tex
%% Autor: Michael Fuhrmann
%% Projekt: Specialist Mathematics Project
%% Thema: Die Whitehead-Asphärizitätsvermutung
%% Sprache: Deutsch
%% Erstellt: 2026-03-12
%% Letzte Änderung: 2026-03-12
%% Build: 164
%%
%% Beschreibung:
%%   Deutschsprachige Version von Paper 78. Die Whitehead-Asphärizitätsvermutung
%%   (1941) besagt: Jeder zusammenhängende Teilkomplex eines asphärischen
%%   2-dimensionalen CW-Komplexes ist asphärisch. Offen seit über 80 Jahren.
%%   CW-Komplexe, Asphärizität, Eilenberg-MacLane-Räume, Cockcroft-Eigenschaft,
%%   partielle Resultate.
%%

\documentclass[12pt,a4paper]{amsart}

%% Pakete für Mathematik, Layout, Encoding und Links
\usepackage{amsmath,amssymb,amsthm,mathtools,enumitem,booktabs,array,hyperref}
\usepackage[utf8]{inputenc}
\usepackage[T1]{fontenc}
\usepackage[ngerman]{babel}
\usepackage{geometry}
\geometry{margin=2.5cm}

%% Theorem-Umgebungen (englisch)
\newtheorem{theorem}{Theorem}[section]
\newtheorem{lemma}[theorem]{Lemma}
\newtheorem{corollary}[theorem]{Korollar}
\newtheorem{proposition}[theorem]{Proposition}

%% Deutsche Theorem-Umgebungen
\newtheorem{satz}[theorem]{Satz}
\newtheorem{korollar}[theorem]{Korollar}
\newtheorem{vermutung}[theorem]{Vermutung}

\theoremstyle{definition}
\newtheorem{definition}[theorem]{Definition}
\newtheorem{definition_de}[theorem]{Definition}
\newtheorem{beispiel}[theorem]{Beispiel}

\theoremstyle{remark}
\newtheorem{bemerkung}[theorem]{Bemerkung}
\newtheorem{remark}[theorem]{Bemerkung}
\newtheorem{conjecture}[theorem]{Vermutung}

%% Mathematische Operatoren und Abkürzungen
\newcommand{\Z}{\mathbb{Z}}
\newcommand{\N}{\mathbb{N}}
\DeclareMathOperator{\cd}{cd}
\DeclareMathOperator{\coker}{coker}
\DeclareMathOperator{\im}{im}

%% Hyperref-Konfiguration
\hypersetup{
    colorlinks=true,
    linkcolor=blue,
    citecolor=blue,
    urlcolor=blue,
    pdftitle={Die Whitehead-Asphärizitätsvermutung},
    pdfauthor={Michael Fuhrmann}
}

%% -------------------------------------------------------------------------
%% Abstract VOR \maketitle
%% -------------------------------------------------------------------------
\begin{abstract}
Die Whitehead-Asphärizitätsvermutung, aufgestellt von J.~H.~C.~Whitehead
im Jahr 1941~\cite{Whitehead1941}, behauptet, dass jeder zusammenhängende
Teilkomplex eines asphärischen 2-dimensionalen CW-Komplexes selbst asphärisch
ist. Trotz ihrer elementaren Formulierung widersteht die Vermutung seit über
80 Jahren jedem Beweis und bleibt eines der zentralen offenen Probleme der
kombinatorischen Gruppentheorie und niedrig-dimensionalen Topologie.

Wir entwickeln die Theorie asphärischer Räume und 2-Komplexe, darunter
Eilenberg--MacLane-Räume, den Satz von van Kampen, den Zusammenhang zwischen
Asphärizität und kohomologischer Dimension, die Cockcroft-Eigenschaft und
den Bezug zur Kohärenz von Gruppen. Wir geben eine Übersicht der wichtigsten
Teilergebnisse: Die Vermutung gilt für Teilkomplexe von 2-Komplexen mit
Einrelatorgruppen (Magnus--Lyndon-Theorie), für 2-Komplexe mit
„diagrammatisch reduzierbaren" Präsentationen (Corson~\cite{Corson1992})
und für viele Klassen von Einrelatorgruppen via den Freiheitssatz.
Der Zusammenhang zum D(2)-Problem von Wall und zur Eilenberg--Ganea-Vermutung
wird ebenfalls behandelt.
\end{abstract}

\title{Die Whitehead-Asphärizitätsvermutung}
\author{Michael Fuhrmann}
\address{Specialist Mathematics Project, Build 165}
\email{specialist-maths@project.local}
\date{2026-03-12}

\maketitle

%% -------------------------------------------------------------------------
\section{Einleitung}
%% -------------------------------------------------------------------------

Im Jahr 1941 fragte J.~H.~C.~Whitehead, ob jeder zusammenhängende
Teilkomplex eines asphärischen 2-dimensionalen CW-Komplexes selbst
asphärisch sein muss~\cite{Whitehead1941}. Ein Raum $X$ ist
\emph{asphärisch}, falls $\pi_n(X) = 0$ für alle $n \geq 2$, d.h.\
falls $X$ ein $K(\pi, 1)$-Raum ist. Für 2-Komplexe ist Asphärizität
äquivalent zum Verschwinden der zweiten Homotopiegruppe $\pi_2(X) = 0$.

Die Bedeutung dieser Frage ergibt sich aus ihren Verbindungen zu:
\begin{itemize}
\item Dem $D(2)$-Problem von C.~T.~C.~Wall~\cite{Wall1965}: ob jeder endliche
      2-dimensionale CW-Komplex, der von einem 3-Komplex dominiert wird,
      homotopieäquivalent zu einem 2-Komplex ist.
\item Kohärenz von Gruppen: eine Gruppe $G$ ist \emph{kohärent}, falls jede
      endlich erzeugte Untergruppe endlich präsentiert ist.
\item Der Eilenberg--Ganea-Vermutung: ob jede Gruppe $G$ der kohomologischen
      Dimension 2 einen 2-dimensionalen $K(G,1)$ besitzt.
\end{itemize}

\begin{vermutung}[Whitehead 1941~\cite{Whitehead1941}]
\label{verm:Whitehead}
Jeder zusammenhängende Teilkomplex eines asphärischen 2-dimensionalen
CW-Komplexes ist asphärisch.
\end{vermutung}

Trotz 80-jähriger Bemühungen ist diese Vermutung vollständig offen.
Es ist kein Gegenbeispiel bekannt, und es existiert kein allgemeiner Beweis.

%% -------------------------------------------------------------------------
\section{CW-Komplexe und die zweite Homotopiegruppe}
%% -------------------------------------------------------------------------

\subsection{CW-Komplexe}

\begin{definition_de}[CW-Komplex]
Ein \emph{CW-Komplex} $X$ ist ein topologischer Raum, der durch sukzessives
Ankleben von Zellen aufgebaut wird: Man beginnt mit einer diskreten Menge
$X^0$ (0-Skelett) und klebt induktiv $n$-Zellen $e_\alpha^n \cong D^n$
via Anheftungsabbildungen $\phi_\alpha \colon \partial D^n \to X^{n-1}$ an.
Das \emph{$n$-Skelett} ist $X^n = X^{n-1} \cup_{\phi_\alpha} e_\alpha^n$.
\end{definition_de}

Ein \emph{2-Komplex} ist ein CW-Komplex der Dimension höchstens 2:
er besteht aus Ecken (0-Zellen), Kanten (1-Zellen) und 2-Zellen (Flächen).

\begin{definition_de}[Asphärischer Raum]
Ein zusammenhängender CW-Komplex $X$ ist \emph{asphärisch}, falls
$\pi_n(X) = 0$ für alle $n \geq 2$.
Äquivalent: Die universelle Überlagerung $\widetilde{X}$ ist kontrahierbar.
Ein asphärischer CW-Komplex $X$ mit $\pi_1(X) = G$ ist ein
\emph{$K(G,1)$-Raum} oder \emph{Eilenberg--MacLane-Raum}.
\end{definition_de}

\subsection{Präsentationskomplexe}

\begin{definition_de}[Präsentationskomplex]
Gegeben einer Gruppenpräsentation $\mathcal{P} = \langle S \mid R \rangle$
ist der \emph{Präsentationskomplex} $X_{\mathcal{P}}$ der 2-Komplex mit
einer 0-Zelle, einer 1-Zelle für jeden Erzeuger $s \in S$ und einer
2-Zelle für jede Relation $r \in R$, angeheftet entsprechend dem Wort $r$.
Dann gilt $\pi_1(X_{\mathcal{P}}) \cong \langle S \mid R \rangle$.
\end{definition_de}

\subsection{Berechnung von $\pi_2$}

\begin{satz}[Hurewicz]
Für einen einfach zusammenhängenden Raum $X$ gilt $\pi_2(X) \cong H_2(X; \Z)$.
Für einen allgemeinen zusammenhängenden Raum $X$ gibt es einen Hurewicz-Homomorphismus
$h \colon \pi_2(X) \to H_2(\widetilde{X}; \Z)$, der surjektiv ist.
\end{satz}

\begin{proposition}
Für einen 2-Komplex $X$ mit Präsentation $\langle S \mid R \rangle$
ist $\pi_2(X)$ ein $\Z[\pi_1(X)]$-Modul und fügt sich in die exakte Sequenz:
\[
    0 \to \pi_2(X) \to C_2(\widetilde{X}) \xrightarrow{\partial_2}
    C_1(\widetilde{X}) \xrightarrow{\partial_1}
    C_0(\widetilde{X}) \to \Z \to 0,
\]
wobei $C_i(\widetilde{X})$ die freien $\Z[\pi_1(X)]$-Module der
$i$-Ketten der universellen Überlagerung $\widetilde{X}$ sind.
\end{proposition}

%% -------------------------------------------------------------------------
\section{Eilenberg--MacLane-Räume und kohomologische Dimension}
%% -------------------------------------------------------------------------

\begin{satz}[Eilenberg--MacLane]
Für jede Gruppe $G$ und jedes $n \geq 1$ existiert ein CW-Komplex $K(G,n)$
(eindeutig bis auf Homotopieäquivalenz) mit $\pi_n(K(G,n)) = G$ und
$\pi_i(K(G,n)) = 0$ für $i \neq n$.
\end{satz}

\begin{definition_de}[Kohomologische Dimension]
Die \emph{kohomologische Dimension} einer Gruppe $G$ ist
\[
    \cd(G) = \sup\{ n \mid H^n(G; M) \neq 0 \text{ für einen $\Z G$-Modul $M$} \}.
\]
\end{definition_de}

\begin{satz}[Stallings 1968, Swan 1969]
$\cd(G) \leq 1$ genau dann, wenn $G$ frei ist.
\end{satz}

\begin{satz}[Eilenberg--Ganea 1957~\cite{EilenbergGanea1957}]
Für $\cd(G) \geq 3$ hat $G$ einen $K(G,1)$ der Dimension $\cd(G)$.
Für $\cd(G) = 2$ hat $G$ einen $K(G,1)$ der Dimension höchstens 3.
\end{satz}

\begin{vermutung}[Eilenberg--Ganea-Vermutung]
\label{verm:EG}
Jede Gruppe $G$ mit $\cd(G) = 2$ hat einen $K(G,1)$ der Dimension 2.
\end{vermutung}

\begin{satz}
Falls die Eilenberg--Ganea-Vermutung~\ref{verm:EG} falsch ist, dann ist
auch die Whitehead-Vermutung~\ref{verm:Whitehead} falsch.
\end{satz}

%% -------------------------------------------------------------------------
\section{Die Cockcroft-Eigenschaft}
%% -------------------------------------------------------------------------

\begin{definition_de}[Cockcroft-Eigenschaft]
Ein 2-Komplex $X$ hat die \emph{Cockcroft-Eigenschaft}, falls jedes Element
von $\pi_2(X)$ unter dem Hurewicz-Homomorphismus auf null in $H_2(X; \Z)$
abgebildet wird.
\end{definition_de}

\begin{satz}[Cockcroft 1954~\cite{Cockcroft1954}]
Ein Präsentationskomplex $X_{\mathcal{P}}$ für $G = \langle S \mid R \rangle$
hat die Cockcroft-Eigenschaft genau dann, wenn $H_2(\widetilde{X}; \Z) = 0$.
\end{satz}

\begin{proposition}
Falls $X$ asphärisch ist, hat $X$ die Cockcroft-Eigenschaft.
\end{proposition}

\begin{proof}
Wenn $X$ asphärisch ist, gilt $\pi_2(X) = 0$, also ist der Hurewicz-Homo\-morphismus
$h \colon \pi_2(X) \to H_2(\widetilde{X})$ trivialerweise die Nullabbildung.
\end{proof}

Die Umkehrung gilt im Allgemeinen nicht. Die Cockcroft-Eigenschaft ist jedoch
eine notwendige, aber nicht hinreichende Bedingung für Asphärizität.

%% -------------------------------------------------------------------------
\section{Einrelatorgruppen und der Freiheitssatz}
%% -------------------------------------------------------------------------

\begin{satz}[Magnus' Freiheitssatz 1930~\cite{Magnus1930}]
Sei $G = \langle x_1, \ldots, x_n \mid r \rangle$ eine Einrelatorgruppe,
wobei $r$ ein zyklisch reduziertes Wort ist, das alle Erzeuger enthält.
Dann ist die Untergruppe, die von einer echten Teilmenge der Erzeuger
$\{x_1, \ldots, x_{n-1}\}$ erzeugt wird, frei.
\end{satz}

\begin{satz}[Lyndon 1950~\cite{Lyndon1950}]
Für eine Einrelatorgruppe $G = \langle S \mid r \rangle$, wobei $r$
keine echte Potenz ist, ist der Präsentationskomplex $X_{\mathcal{P}}$
asphärisch.
\end{satz}

\begin{proof}[Beweisskizze]
Lyndon's Identitätssatz zeigt, dass $\pi_2(X_{\mathcal{P}})$ als
$\Z G$-Modul von „sphärischen Diagrammen" erzeugt wird, die er als
trivial nachweist, wenn $r$ keine echte Potenz ist. Der Beweis verwendet
die Magnus-Methode (Induktion über die Erzeugeranzahl) und den Freiheitssatz.
\end{proof}

\begin{satz}[Whitehead-Vermutung für Einrelator-Komplexe]
Sei $X$ ein asphärischer Einrelator-Präsentationskomplex (Lyndon-Bedingung
erfüllt). Dann ist jeder zusammenhängende Teilkomplex $Y \subseteq X$
asphärisch.
\end{satz}

\begin{proof}
Jeder zusammenhängende Teilkomplex eines Einrelator-Präsentationskomplexes
ist entweder das 1-Skelett (ein Kreis-Wedge, asphärisch) oder der gesamte
Komplex (asphärisch nach Hypothese).
\end{proof}

%% -------------------------------------------------------------------------
\section{Diagrammatische Reduzierbarkeit}
%% -------------------------------------------------------------------------

\begin{definition_de}[Diagrammatisch reduzierbare Präsentation]
Eine Präsentation $\mathcal{P} = \langle S \mid R \rangle$ ist
\emph{diagrammatisch reduzierbar} (DR), falls jedes sphärische Diagramm
über $\mathcal{P}$ (d.h.\ jede zelluläre Abbildung $f \colon S^2 \to X_{\mathcal{P}}$)
durch eine Folge von „Faltungsoperationen" zur konstanten Abbildung
reduziert werden kann.
\end{definition_de}

\begin{satz}[Sieradski 1983~\cite{Sieradski1983}]
Wenn $\mathcal{P}$ diagrammatisch reduzierbar ist, dann ist
$X_{\mathcal{P}}$ asphärisch.
\end{satz}

\begin{satz}[Corson 1992~\cite{Corson1992}]
\label{satz:Corson}
Jeder zusammenhängende Teilkomplex eines diagrammatisch reduzierbaren
2-Komplexes ist asphärisch.
\end{satz}

\begin{proof}[Beweisskizze]
Corson zeigt, dass die DR-Eigenschaft „erblich" ist in folgendem Sinne:
Wenn $\mathcal{P}$ DR ist und $\mathcal{Q}$ durch Entfernen einiger
Relationen aus $\mathcal{P}$ entsteht, dann ist $\mathcal{Q}$ ebenfalls DR.
Die diagrammatische Reduzierbarkeit der Präsentation eines Teilkomplexes
folgt dann aus der DR-Eigenschaft des Gesamtkomplexes, und DR impliziert
Asphärizität.
\end{proof}

%% -------------------------------------------------------------------------
\section{Kohärenz von Gruppen}
%% -------------------------------------------------------------------------

\begin{definition_de}[Kohärente Gruppe]
Eine Gruppe $G$ ist \emph{kohärent}, falls jede endlich erzeugte Untergruppe
von $G$ endlich präsentiert ist.
\end{definition_de}

\begin{satz}[Scott 1973~\cite{Scott1973}]
Fundamentalgruppen kompakter 3-Mannigfaltigkeiten sind kohärent.
\end{satz}

\begin{satz}[Stallings 1968]
Freie Gruppen sind kohärent.
\end{satz}

\begin{bemerkung}
Ein Ansatz zur Whitehead-Vermutung über Kohärenz: Falls $X$ ein asphärischer
2-Komplex mit $\pi_1(X) = G$ ist und $Y \subseteq X$ ein zusammenhängender
Teilkomplex mit $\pi_1(Y) = H \leq G$ ist, dann ist die Asphärizität von $Y$
äquivalent dazu, dass $\cd(H) \leq 1$ (d.h.\ $H$ frei ist) oder dass $H$
einen 2-dimensionalen $K(H,1)$ besitzt. Die Kohärenz von $G$ stellt sicher,
dass $H$ endlich präsentiert ist, liefert aber direkt keine Asphärizität.
\end{bemerkung}

%% -------------------------------------------------------------------------
\section{Das D(2)-Problem von Wall}
%% -------------------------------------------------------------------------

\begin{satz}[Wall 1965~\cite{Wall1965}]
Jeder zusammenhängende endliche CW-Komplex $X$ mit $H^n(X; M) = 0$
für alle $n \geq 3$ und alle $\Z[\pi_1(X)]$-Moduln $M$ ist
homotopieäquivalent zu einem 3-dimensionalen CW-Komplex.
\end{satz}

Das $D(2)$-Problem fragt, ob solch ein Raum homotopieäquivalent zu einem
2-dimensionalen Komplex ist:

\begin{vermutung}[$D(2)$-Problem, Wall~\cite{Wall1965}]
\label{verm:D2}
Falls $X$ ein endlicher zusammenhängender CW-Komplex ist, der von einem
2-Komplex dominiert wird, und $H^n(X; M) = 0$ für $n \geq 3$ und alle
$\Z[\pi_1 X]$-Moduln $M$, dann ist $X$ homotopieäquivalent zu einem
endlichen 2-Komplex.
\end{vermutung}

\begin{satz}[Johnson 2003~\cite{Johnson2003}]
Das $D(2)$-Problem ist äquivalent zur Whitehead-Vermutung für endliche
2-Komplexe.
\end{satz}

Diese Äquivalenz ist überraschend und tief: Sie verbindet ein Problem
über Teilraumstruktur (Whitehead) mit einem Problem über Dimensionsreduktion
(Wall).

%% -------------------------------------------------------------------------
\section{Teilergebnisse und Spezialfälle}
%% -------------------------------------------------------------------------

\begin{satz}[Magnus-Lyndon, Einrelatorgruppen]
Die Whitehead-Vermutung gilt für Teilkomplexe von Präsentationskomplexen
von Einrelatorgruppen mit torsionsfreiem Relator.
\end{satz}

\begin{satz}[Corson 1992, Satz~\ref{satz:Corson}]
Die Whitehead-Vermutung gilt für Teilkomplexe diagrammatisch
reduzierbarer 2-Komplexe.
\end{satz}

\begin{satz}[Howie 1979~\cite{Howie1979}]
Die Whitehead-Vermutung gilt, wenn der Teilkomplex einfach zusammenhängend ist.
\end{satz}

\begin{satz}[Asphärische Gruppen der kohomologischen Dimension 2]
Für Gruppen $G$ der kohomologischen Dimension 2 (z.B.\ Flächengruppen,
Torusknoten-Gruppen) hat jede Untergruppe kohomologische Dimension $\leq 2$,
was Evidenz für Whitehead's Vermutung in diesen Fällen liefert.
\end{satz}

%% -------------------------------------------------------------------------
\section{Gegenbeispielversuche und Hindernisse}
%% -------------------------------------------------------------------------

Die Schwierigkeit der Whitehead-Vermutung zeigt sich darin, dass viele
potenzielle Gegenbeispielkonstruktionen aus subtilen Gründen scheitern.

\begin{bemerkung}
Das einfachste potenzielle Gegenbeispiel wäre: ein asphärischer 2-Komplex $X$
und ein Teilkomplex $Y \subseteq X$ mit $\pi_2(Y) \neq 0$. Für ein solches
$Y$ wäre $\pi_2(Y)$ ein nichttrivialer $\Z[\pi_1(Y)]$-Modul. Jedes
Element $\alpha \in \pi_2(Y)$ bildet sich aber auf ein Element von
$\pi_2(X) = 0$ ab (da $X$ asphärisch ist), also ist $\alpha$ schon
in $X$ nullhomotop. Die Whitehead-Vermutung besagt, dass es bereits
in $Y$ nullhomotop ist.
\end{bemerkung}

\begin{lemma}[Nicht-erbliche Asphärizität]
Asphärizität ist im Allgemeinen nicht erblich unter Teilraumeinbettungen.
Beispielsweise ist der Torus $T^2 = K(\Z^2, 1)$ asphärisch, aber wenn
man eine 2-Zelle entfernt, erhält man das 1-Skelett $S^1 \vee S^1$,
das ebenfalls asphärisch ist. Bei komplizierteren Beispielen können
Teilkomplexe jedoch nichttriviales $\pi_2$ haben.
\end{lemma}

%% -------------------------------------------------------------------------
\section{Offene Probleme}
%% -------------------------------------------------------------------------

\begin{vermutung}[Whitehead, Verm.~\ref{verm:Whitehead}]
Jeder zusammenhängende Teilkomplex eines asphärischen 2-Komplexes
ist asphärisch.
\end{vermutung}

\begin{vermutung}[Eilenberg--Ganea, Verm.~\ref{verm:EG}]
Jede Gruppe $G$ mit $\cd(G) = 2$ hat einen $K(G,1)$ der Dimension 2.
\end{vermutung}

\begin{vermutung}[$D(2)$-Problem, Verm.~\ref{verm:D2}]
Jeder endliche 2-Komplex, der von einem 3-Komplex dominiert wird,
ist homotopieäquivalent zu einem 2-Komplex.
\end{vermutung}

\begin{vermutung}[Asphärizität aller Einrelatorgruppen]
\label{verm:einrelator-asph}
Jede Einrelatorgruppe $G = \langle S \mid r \rangle$ (ohne Torsion)
hat einen asphärischen Präsentationskomplex.
\end{vermutung}

\begin{vermutung}[Kohärenzvermutung für hyperbolische Gruppen]
Jede Einrelatorgruppe ist kohärent.
\end{vermutung}

\begin{bemerkung}
Die Vermutungen~\ref{verm:Whitehead}, \ref{verm:EG} und \ref{verm:D2} sind
miteinander verbunden: Ein Gegenbeispiel zur Eilenberg--Ganea-Vermutung würde
ein Gegenbeispiel zur Whitehead-Vermutung liefern; das $D(2)$-Problem ist
äquivalent zur Whitehead-Vermutung für endliche Komplexe (Johnson 2003).
Ein gemeinsamer Ansatz zur Lösung aller drei wäre eine neue Konstruktion
asphärischer Komplexe via CAT(0)-Geometrie oder hyperbolischer Gruppentheorie.
\end{bemerkung}

%% -------------------------------------------------------------------------
\section{Zusammenfassung}
%% -------------------------------------------------------------------------

Die Whitehead-Asphärizitätsvermutung berührt trotz ihrer elementaren
Formulierung tiefe Aspekte der kombinatorischen Gruppentheorie, algebraischen
Topologie und 4-Mannigfaltigkeiten-Topologie. Die Verbindung zum $D(2)$-Problem
von Wall zeigt, dass sie keine isolierte Kuriosität ist, sondern mit
grundlegenden Fragen über die Existenz und Klassifikation 2-dimensionaler
klassifizierender Räume verwoben ist. Teilergebnisse (Corson's Satz über
diagrammatisch reduzierbare Komplexe, der Magnus-Lyndon-Satz für
Einrelatorgruppen, Howie's Resultat für einfach zusammenhängende Teilkomplexe)
bestätigen die Vermutung in wichtigen Spezialfällen, aber ein allgemeiner
Beweis bleibt außer Reichweite.

\begin{thebibliography}{99}

\bibitem{Whitehead1941}
J.~H.~C.~Whitehead,
\textit{On adding relations to homotopy groups},
Ann.\ of Math.\ \textbf{42} (1941), 409--428.

\bibitem{Wall1965}
C.~T.~C.~Wall,
\textit{Finiteness conditions for CW-complexes},
Ann.\ of Math.\ \textbf{81} (1965), 56--69.

\bibitem{EilenbergGanea1957}
S.~Eilenberg, T.~Ganea,
\textit{On the Lusternik-Schnirelmann category of abstract groups},
Ann.\ of Math.\ \textbf{65} (1957), 517--518.

\bibitem{Magnus1930}
W.~Magnus,
\textit{Über diskontinuierliche Gruppen mit einer definierenden Relation
(der Freiheitssatz)},
J.\ Reine Angew.\ Math.\ \textbf{163} (1930), 141--165.

\bibitem{Lyndon1950}
R.~C.~Lyndon,
\textit{Cohomology theory of groups with a single defining relation},
Ann.\ of Math.\ \textbf{52} (1950), 650--665.

\bibitem{Cockcroft1954}
W.~H.~Cockcroft,
\textit{On two-dimensional aspherical complexes},
Proc.\ London Math.\ Soc.\ \textbf{4} (1954), 375--384.

\bibitem{Sieradski1983}
A.~J.~Sieradski,
\textit{A coloring invariant for topological graph theory},
J.\ Combin.\ Theory Ser.\ B \textbf{34} (1983), 241--246.

\bibitem{Corson1992}
J.~M.~Corson,
\textit{Conformally aspherical 2-complexes},
J.\ Pure Appl.\ Algebra \textbf{78} (1992), 11--22.

\bibitem{Scott1973}
G.~P.~Scott,
\textit{Finitely generated 3-manifold groups are finitely presented},
J.\ London Math.\ Soc.\ \textbf{6} (1973), 437--440.

\bibitem{BridsonHaefliger1999}
M.~R.~Bridson, A.~Haefliger,
\textit{Metric Spaces of Non-Positive Curvature},
Grundlehren der mathematischen Wissenschaften \textbf{319},
Springer-Verlag, Berlin, 1999.

\bibitem{Howie1979}
J.~Howie,
\textit{Aspherical and acyclic 2-complexes},
J.\ London Math.\ Soc.\ \textbf{20} (1979), 549--558.

\bibitem{Johnson2003}
F.~E.~A.~Johnson,
\textit{Stable Modules and the $D(2)$-Problem},
London Mathematical Society Lecture Note Series \textbf{301},
Cambridge University Press, 2003.

\end{thebibliography}

\end{document}
